%% ==================================================================
%% TOAN12_C5_BAI14.tex -- Toán 12, Chương V, Bài 14
%% Cấu trúc: Lí thuyết trọng tâm -> Ví dụ minh hoạ (chừa chỗ trống) -> Bài tập luyện tập
%% ==================================================================
\baihoc{14}{Phương trình mặt phẳng}

\lithuyet

\begin{dinhnghia}[Vectơ pháp tuyến của mặt phẳng]
\begin{minipage}[c]{.62\linewidth}
Vectơ $\vec n\ne\vec0$ được gọi là \textbf{vectơ pháp tuyến} của mặt phẳng $(\alpha)$ nếu giá của $\vec n$ vuông góc với $(\alpha)$.

\smallskip\textbf{\color{dcnhan}Chú ý.}
\begin{itemize}
  \item Mặt phẳng hoàn toàn xác định khi biết một điểm và một vectơ pháp tuyến của nó.
  \item Nếu $\vec n$ là một vectơ pháp tuyến của mặt phẳng $(\alpha)$ thì $k\vec n$ (với $k$ là một số khác $0$) cũng là một vectơ pháp tuyến của $(\alpha)$.
\end{itemize}
\end{minipage}\hfill
\begin{minipage}[c]{.35\linewidth}\centering
\begin{tikzpicture}[scale=.9,font=\small]
  \fill[dcnen,draw=dcdam,line width=.7pt] (0,0)--(3.2,0)--(4,1.2)--(.8,1.2)--cycle;
  \node at (.55,.25){$(\alpha)$};
  \draw[vecto,dcnhan] (2,.6)--(2,2.2) node[right]{$\vec n$};
  \fill[dcdam] (2,.6) circle (1.3pt);
  \draw[dcdam] (2,.78)--(2.18,.78)--(2.18,.6);
\end{tikzpicture}
\end{minipage}
\end{dinhnghia}

\begin{dinhnghia}[Tích có hướng của hai vectơ]
\begin{minipage}[c]{.62\linewidth}
Trong không gian $Oxyz$, cho hai vectơ $\vec u=(a;b;c)$ và $\vec v=(a';b';c')$. Khi đó vectơ
\[\vec n=(bc'-b'c;\ ca'-c'a;\ ab'-a'b)\]
vuông góc với cả hai vectơ $\vec u$ và $\vec v$, được gọi là \textbf{tích có hướng} của $\vec u$ và $\vec v$, kí hiệu là $[\vec u,\vec v]$.
\end{minipage}\hfill
\begin{minipage}[c]{.35\linewidth}\centering
\begin{tikzpicture}[scale=.9,font=\small]
  \draw[vecto,dcchinh] (0,0)--(-1,-.8) node[left]{$\vec u$};
  \draw[vecto,dcchinh] (0,0)--(1.6,0) node[right]{$\vec v$};
  \draw[vecto,dcnhan] (0,0)--(0,1.6) node[above]{$[\vec u,\vec v]$};
\end{tikzpicture}
\end{minipage}
\end{dinhnghia}

\begin{chuy}
\begin{itemize}
  \item $[\vec u,\vec v]=\vec0$ khi và chỉ khi $\vec u$, $\vec v$ cùng phương.
  \item Với bốn số $x,y,x',y'$, ta kí hiệu $\begin{vmatrix}x&y\\x'&y'\end{vmatrix}=xy'-x'y$. Khi đó tích có hướng của $\vec u=(a;b;c)$ và $\vec v=(a';b';c')$ là
  \[[\vec u,\vec v]=\left(\begin{vmatrix}b&c\\b'&c'\end{vmatrix};\ \begin{vmatrix}c&a\\c'&a'\end{vmatrix};\ \begin{vmatrix}a&b\\a'&b'\end{vmatrix}\right).\]
\end{itemize}
\end{chuy}

\begin{dinhnghia}[Cặp vectơ chỉ phương của mặt phẳng]
\begin{itemize}
  \item Trong không gian $Oxyz$, hai vectơ $\vec u$, $\vec v$ được gọi là \textbf{cặp vectơ chỉ phương} của mặt phẳng $(P)$ nếu chúng không cùng phương và có giá nằm trong hoặc song song với mặt phẳng $(P)$.
  \item Nếu $\vec u$, $\vec v$ là cặp vectơ chỉ phương của $(P)$ thì $[\vec u,\vec v]$ là một vectơ pháp tuyến của $(P)$.
\end{itemize}
\end{dinhnghia}

\begin{dinhnghia}[Phương trình tổng quát của mặt phẳng]
Trong không gian $Oxyz$, mỗi mặt phẳng đều có phương trình dạng $Ax+By+Cz+D=0$, trong đó $A$, $B$, $C$ không đồng thời bằng $0$, được gọi là \textbf{phương trình tổng quát} của mặt phẳng đó.

\smallskip\textbf{\color{dcnhan}Chú ý.} Trong không gian $Oxyz$, mỗi phương trình $Ax+By+Cz+D=0$ (các hệ số $A$, $B$, $C$ không đồng thời bằng $0$) xác định một mặt phẳng nhận $\vec n=(A;B;C)$ làm một vectơ pháp tuyến.
\end{dinhnghia}

\begin{congthuc}[Lập phương trình tổng quát của mặt phẳng]
\begin{itemize}
  \item \textbf{Đi qua một điểm và biết vectơ pháp tuyến.} Nếu mặt phẳng $(\alpha)$ đi qua điểm $M_0(x_0;y_0;z_0)$ và có vectơ pháp tuyến $\vec n=(A;B;C)$ thì có phương trình là:
  \[A(x-x_0)+B(y-y_0)+C(z-z_0)=0\Leftrightarrow Ax+By+Cz+D=0,\ \text{với } D=-(Ax_0+By_0+Cz_0).\]
  \item \textbf{Đi qua điểm $M$ và biết cặp vectơ chỉ phương $\vec u$, $\vec v$:} tìm vectơ pháp tuyến $\vec n=[\vec u,\vec v]$, rồi lập phương trình tổng quát của mặt phẳng đi qua $M$ và có vectơ pháp tuyến $\vec n$.
  \item \textbf{Đi qua ba điểm không thẳng hàng $A$, $B$, $C$:} tìm cặp vectơ chỉ phương $\vt{AB}$, $\vt{AC}$; tìm vectơ pháp tuyến $\vec n=\big[\vt{AB},\vt{AC}\big]$; lập phương trình tổng quát của mặt phẳng đi qua $A$ và có vectơ pháp tuyến $\vec n$.
\end{itemize}
\end{congthuc}

\begin{kienthuc}[Phương trình mặt phẳng theo đoạn chắn]
\begin{minipage}[c]{.6\linewidth}
Mặt phẳng $(\alpha)$ không đi qua gốc toạ độ và cắt ba trục $Ox$, $Oy$, $Oz$ tương ứng tại các điểm $A(a;0;0)$, $B(0;b;0)$, $C(0;0;c)$ ($a,b,c\ne0$) có phương trình
\[\frac xa+\frac yb+\frac zc=1.\]
\end{minipage}\hfill
\begin{minipage}[c]{.37\linewidth}\centering
\begin{tikzpicture}[scale=.85,font=\scriptsize,x={(-0.5cm,-0.45cm)},y={(1cm,0cm)},z={(0cm,1cm)}]
  \fill[dcchinh!15] (2.2,0,0)--(0,2.4,0)--(0,0,2)--cycle;
  \draw[truc] (0,0,0)--(3,0,0) node[below left]{$x$}; \draw[truc] (0,0,0)--(0,3.2,0) node[right]{$y$}; \draw[truc] (0,0,0)--(0,0,2.6) node[above]{$z$};
  \draw[dcdam,line width=.8pt] (2.2,0,0)--(0,2.4,0)--(0,0,2)--cycle;
  \node[left] at (2.2,0,0){$A(a;0;0)$}; \node[below right] at (0,2.4,0){$B(0;b;0)$}; \node[right] at (0,0,2){$C(0;0;c)$};
  \node[above left] at (0,0,0){$O$};
\end{tikzpicture}
\end{minipage}
\end{kienthuc}

\begin{congthuc}[Điều kiện để hai mặt phẳng vuông góc, song song]
Trong không gian $Oxyz$, cho hai mặt phẳng $(\alpha)\colon Ax+By+Cz+D=0$, $(\beta)\colon A'x+B'y+C'z+D'=0$, với hai vectơ pháp tuyến $\vec n=(A;B;C)$, $\vec n'=(A';B';C')$ tương ứng. Khi đó:
\begin{itemize}
  \item $(\alpha)\perp(\beta)\Leftrightarrow\vec n\perp\vec n'\Leftrightarrow AA'+BB'+CC'=0$.
  \item $(\alpha)\parallel(\beta)\Leftrightarrow\begin{cases}\vec n'=k\vec n\\D'\ne kD\end{cases}$ với $k$ nào đó.
\end{itemize}
\hinh{\begin{tikzpicture}[scale=.75,font=\small]
  \begin{scope}
    \fill[dcnen,draw=dcdam] (0,0)--(3,0)--(3.8,1)--(.8,1)--cycle;
    \fill[dcvang,draw=dcdam,opacity=.9] (.8,1)--(0,0)--(0,2)--(.8,3)--cycle;
    \draw[vecto,dcnhan] (2.3,.5)--(2.3,1.9) node[right]{$\vec n'$};
    \draw[vecto,dcchinh] (.5,1.9)--(1.6,1.9) node[right]{$\vec n$};
    \node at (1.9,-.5){$(\alpha)\perp(\beta)$};
  \end{scope}
  \begin{scope}[xshift=6.2cm]
    \fill[dcnen,draw=dcdam] (0,0)--(3,0)--(3.8,1)--(.8,1)--cycle;
    \fill[dcnen,draw=dcdam] (0,1.6)--(3,1.6)--(3.8,2.6)--(.8,2.6)--cycle;
    \draw[vecto,dcchinh] (2,.5)--(2,1.25) node[right]{$\vec n'$};
    \draw[vecto,dcchinh] (2,2.1)--(2,2.9) node[right]{$\vec n$};
    \node at (1.9,-.5){$(\alpha)\parallel(\beta)$};
  \end{scope}
\end{tikzpicture}}
\end{congthuc}

\begin{chuy}
\begin{itemize}
  \item Nếu hai mặt phẳng vuông góc với nhau thì vectơ pháp tuyến của mặt phẳng này có giá song song hoặc nằm trong mặt phẳng kia.
  \item Nếu hai mặt phẳng song song với nhau thì vectơ pháp tuyến của mặt phẳng này cũng là vectơ pháp tuyến của mặt phẳng kia.
  \item Hai mặt phẳng $(\alpha)$ và $(\beta)$ trùng nhau khi và chỉ khi tồn tại số $k$ khác $0$ sao cho $A'=kA$, $B'=kB$, $C'=kC$, $D'=kD$.
\end{itemize}
\end{chuy}

\begin{congthuc}[Khoảng cách từ một điểm đến một mặt phẳng]
\begin{minipage}[c]{.62\linewidth}
Trong không gian $Oxyz$, khoảng cách từ điểm $M(x_0;y_0;z_0)$ đến mặt phẳng $(P)\colon Ax+By+Cz+D=0$ là:
\[\mathrm d\big(M,(P)\big)=\frac{|Ax_0+By_0+Cz_0+D|}{\sqrt{A^2+B^2+C^2}}.\]
\end{minipage}\hfill
\begin{minipage}[c]{.35\linewidth}\centering
\begin{tikzpicture}[scale=.85,font=\small]
  \fill[dcnen,draw=dcdam,line width=.7pt] (0,0)--(3.2,0)--(4,1.2)--(.8,1.2)--cycle;
  \node at (.6,.25){$(P)$};
  \draw[dcnhan,line width=1pt] (2.2,.6)--(2.2,2.2);
  \fill[dcdam] (2.2,2.2) circle (1.5pt) node[right]{$M$}; \fill[dcdam] (2.2,.6) circle (1.5pt) node[below right]{$N$};
  \draw[dcdam] (2.2,.78)--(2.38,.78)--(2.38,.6);
\end{tikzpicture}
\end{minipage}
\end{congthuc}

%% ------------------------------------------------ VÍ DỤ
\viduminhhoa
\subsubsection{Vectơ pháp tuyến, cặp vectơ chỉ phương, tích có hướng}
\begin{vidu}
Trong không gian $Oxyz$, cho $\vec u=(1;-2;0)$ và $\vec v=(3;1;-4)$. Tính $[\vec u,\vec v]$.
\begin{loigiai}
$[\vec u,\vec v]=\left(\begin{vmatrix}-2&0\\1&-4\end{vmatrix};\begin{vmatrix}0&1\\-4&3\end{vmatrix};\begin{vmatrix}1&-2\\3&1\end{vmatrix}\right)=(8;4;7)$.
\end{loigiai}
\end{vidu}

\begin{vidu}
Trong không gian $Oxyz$, cho mặt phẳng $(\alpha)\colon x+2y-z+1=0$. a) Chỉ ra một vectơ pháp tuyến của $(\alpha)$. b) Vectơ $\vec m=(2;4;-2)$ có là vectơ pháp tuyến của $(\alpha)$ không? c) Trong hai điểm $A(1;3;2)$, $B(1;1;4)$, điểm nào thuộc $(\alpha)$?
\begin{loigiai}
a) $\vec n=(1;2;-1)$. b) Có, vì $\vec m=2\vec n$. c) $1+6-2+1\ne0$, $1+2-4+1=0$: chỉ $B$ thuộc $(\alpha)$.
\end{loigiai}
\end{vidu}

\subsubsection{Lập phương trình mặt phẳng}
\begin{vidu}
Trong không gian $Oxyz$, viết phương trình mặt phẳng $(\alpha)$: a) đi qua $M(2;-1;0)$ và có vectơ pháp tuyến $\vec n=(3;-4;6)$; b) đi qua $M(1;2;-4)$ và vuông góc với trục $Oz$.
\begin{loigiai}
a) $3(x-2)-4(y+1)+6z=0\Leftrightarrow3x-4y+6z-10=0$. b) Vectơ pháp tuyến $\vec k=(0;0;1)$: $z+4=0$.
\end{loigiai}
\end{vidu}

\begin{vidu}
Trong không gian $Oxyz$, cho ba điểm $A(2;1;-1)$, $B(3;2;1)$, $C(3;1;4)$. a) Chứng minh $A$, $B$, $C$ không thẳng hàng. b) Viết phương trình mặt phẳng $(ABC)$.
\begin{loigiai}
a) $\vt{AB}=(1;1;2)$, $\vt{AC}=(1;0;5)$ không cùng phương. b) $\vec n=\big[\vt{AB},\vt{AC}\big]=(5;-3;-1)$; $(ABC)\colon5(x-2)-3(y-1)-(z+1)=0\Leftrightarrow5x-3y-z-8=0$.
\end{loigiai}
\end{vidu}

\begin{vidu}
Trong không gian $Oxyz$, viết phương trình mặt phẳng $(P)$ đi qua hai điểm $A(1;2;-2)$, $B(2;4;1)$ và vuông góc với mặt phẳng $(Q)\colon x+3y+z-1=0$.
\begin{loigiai}
$(P)$ có cặp vectơ chỉ phương $\vt{AB}=(1;2;3)$ và $\vec n_Q=(1;3;1)$, nên $\vec n_P=\big[\vt{AB},\vec n_Q\big]=(-7;2;1)$. $(P)\colon-7(x-1)+2(y-2)+(z+2)=0\Leftrightarrow7x-2y-z-5=0$.
\end{loigiai}
\end{vidu}

\begin{vidu}
Trong không gian $Oxyz$, viết phương trình mặt phẳng đi qua ba điểm $A(2;0;0)$, $B(0;3;0)$, $C(0;0;-1)$.
\begin{loigiai}
Phương trình theo đoạn chắn: $\dfrac x2+\dfrac y3+\dfrac z{-1}=1\Leftrightarrow3x+2y-6z-6=0$.
\end{loigiai}
\end{vidu}

\subsubsection{Vị trí tương đối của hai mặt phẳng}
\begin{vidu}
Trong không gian $Oxyz$, xét vị trí tương đối của các cặp mặt phẳng: a) $(\alpha)\colon x-3y+2z+1=0$ và $(\beta)\colon5x+y-z+2=0$; b) $(P)\colon x+2y-2z+1=0$ và $(Q)\colon2x+4y-4z+5=0$.
\begin{loigiai}
a) $1\cdot5+(-3)\cdot1+2\cdot(-1)=0$ nên $(\alpha)\perp(\beta)$. b) $\vec n_Q=2\vec n_P$ và $5\ne2\cdot1$ nên $(P)\parallel(Q)$.
\end{loigiai}
\end{vidu}

\subsubsection{Khoảng cách từ một điểm đến một mặt phẳng}
\begin{vidu}
Trong không gian $Oxyz$: a) Tính khoảng cách từ $M(1;2;-1)$ đến $(P)\colon x+2y-2z+5=0$. b) Tính khoảng cách giữa hai mặt phẳng song song $(P)\colon x+y+z+2=0$ và $(Q)\colon x+y+z+6=0$.
\begin{loigiai}
a) $\mathrm d(M,(P))=\dfrac{|1+4+2+5|}{\sqrt{1+4+4}}=4$. b) Lấy $A(-2;0;0)\in(P)$: $\mathrm d((P),(Q))=\mathrm d(A,(Q))=\dfrac{|-2+6|}{\sqrt3}=\dfrac{4\sqrt3}3$.
\end{loigiai}
\end{vidu}

\begin{vidu}
Góc quan sát ngang của một camera là $115^\circ$. Trong không gian $Oxyz$, camera đặt tại điểm $C(1;2;4)$ và chiếu thẳng về phía mặt phẳng $(P)\colon x+2y+2z+3=0$. Hỏi vùng quan sát được trên $(P)$ là hình tròn có bán kính bằng bao nhiêu (làm tròn đến hàng phần mười)?
\hinh{\begin{tikzpicture}[scale=.6,font=\scriptsize]
  \fill[dcnen,draw=dcdam] (-2.6,0)--(2.2,0)--(3,.9)--(-1.8,.9)--cycle;
  \draw[dcchinh,dashed] (.2,.45) ellipse (1.6 and .3);
  \draw[dcdam] (.2,3)--(-1.4,.45) (.2,3)--(1.8,.45); \draw[dcnhan,line width=.8pt] (.2,3)--(.2,.45);
  \fill (.2,3) circle (1.5pt) node[above]{$C$};
\end{tikzpicture}}
\begin{loigiai}
Khoảng cách từ $C$ đến $(P)$: $h=\dfrac{|1+4+8+3|}{3}=\dfrac{16}3$. Bán kính $r=h\cdot\tan57{,}5^\circ\approx8{,}4$.
\end{loigiai}
\end{vidu}

\begin{viduds}
Trong không gian $Oxyz$, cho mặt phẳng $(P)\colon2x-y+2z-3=0$ và điểm $M(1;1;1)$. Xét tính đúng sai của các mệnh đề sau:
\begin{dungsai}
  \y Vectơ $\vec n=(2;-1;2)$ là một vectơ pháp tuyến của $(P)$. \oDS
  \y Điểm $M$ thuộc mặt phẳng $(P)$. \oDS
  \y Khoảng cách từ gốc toạ độ $O$ đến $(P)$ bằng $3$. \oDS
  \y Mặt phẳng $(Q)\colon2x-y+2z+3=0$ song song với $(P)$ và khoảng cách giữa chúng bằng $2$. \oDS
\end{dungsai}
\begin{loigiai}
\tuduy Đọc vectơ pháp tuyến; thay toạ độ; áp dụng công thức khoảng cách; hai mặt phẳng song song có cùng vectơ pháp tuyến.

a) \textbf{Đúng}. b) \textbf{Đúng}: $2-1+2-3=0$. c) \textbf{Sai}: $\mathrm d(O,(P))=\dfrac33=1$. d) \textbf{Đúng}: $M\in(P)$, $\mathrm d(M,(Q))=\dfrac{|2-1+2+3|}3=2$.
\end{loigiai}
\end{viduds}
\subsubsection{Bài toán thực tế}
\begin{vidu}
Trong một hệ trục $Oxyz$ (đơn vị: mét), một tấm kính phẳng được đặt đi qua ba điểm $A(4;0;0)$, $B(0;3;0)$, $C(0;0;4)$. Một đèn nhỏ đặt tại điểm $E(2;1;3)$. Khoảng cách (mét) từ đèn đến tấm kính bằng bao nhiêu (làm tròn kết quả đến hàng phần trăm)?
\hinh{\begin{tikzpicture}[scale=0.85,font=\footnotesize,line join=round,line cap=round,>=Stealth]
\draw[->,gray] (0.000,0.000)--(-1.260,-0.945) node[below left]{$x$};
\draw[->,gray] (0.000,0.000)--(4.300,0.000) node[right]{$y$};
\draw[->,gray] (0.000,0.000)--(0.000,3.200) node[above]{$z$};
\node[below right,gray] at (0.000,0.000){$O$};
\fill[gray!25,opacity=.8] (-0.936,-0.702)--(3.400,0.000)--(0.000,2.500)--cycle; \draw[thick] (-0.936,-0.702) node[left]{$A$}--(3.400,0.000) node[below]{$B$}--(0.000,2.500) node[left]{$C$}--cycle;
\shade[ball color=gray!20] (0.640,0.030) circle (.1); \node[right] at (0.640,0.030){$E$};
\end{tikzpicture}}

\begin{loigiai}
\giai Mặt phẳng $(ABC)\colon \dfrac{x}{4}+\dfrac{y}{3}+\dfrac{z}{4}=1\Leftrightarrow 3x+4y+3z-12=0$. $d=\dfrac{|28|}{\sqrt{544}}\approx 1{,}2$. Đáp số: $1{,}2$.
\end{loigiai}
\end{vidu}
\begin{vidu}
Trong không gian $Oxyz$ (đơn vị: mét), mặt một sườn đồi được mô hình hoá bởi mặt phẳng $(P)\colon 12x-3y+4z-7=0$. Một flycam đang ở vị trí $F(-2;-4;11)$. Khoảng cách (mét) từ flycam đến mặt sườn đồi bằng bao nhiêu (làm tròn kết quả đến hàng phần trăm)?
\hinh{\begin{tikzpicture}[scale=0.85,font=\footnotesize,line join=round,line cap=round,>=Stealth]
\draw[->,gray] (0.000,0.000)--(-1.440,-1.080) node[below left]{$x$};
\draw[->,gray] (0.000,0.000)--(5.000,0.000) node[right]{$y$};
\draw[->,gray] (0.000,0.000)--(0.000,3.300) node[above]{$z$};
\node[below right,gray] at (0.000,0.000){$O$};
\fill[gray!10] (-0.500,0.000)--(-1.760,-0.945)--(3.240,-0.945)--(4.500,0.000)--cycle;
\draw[dashed] (0.300,1.525)--(0.300,-0.675); \fill (0.300,-0.675) circle (.04);
\filldraw[fill=gray!40] (0.300,1.525) ++(-.18,0) rectangle ++(.36,.1); \node[above] at (0.300,1.525){$A$};
\draw[dashed] (3.240,2.530)--(3.240,-0.270); \fill (3.240,-0.270) circle (.04);
\filldraw[fill=gray!40] (3.240,2.530) ++(-.18,0) rectangle ++(.36,.1); \node[above] at (3.240,2.530){$B$};
\draw[thick,dotted] (0.300,1.525)--(3.240,2.530);
\end{tikzpicture}}

\begin{loigiai}
\giai $d(F,(P))=\dfrac{|25|}{13}=\dfrac{25}{13}$ (m). Đáp số: $1{,}92$.
\end{loigiai}
\end{vidu}
\begin{vidu}
Trong một hệ trục $Oxyz$ (đơn vị: mét), một tấm kính phẳng được đặt đi qua ba điểm $A(5;0;0)$, $B(0;4;0)$, $C(0;0;4)$. Một đèn nhỏ đặt tại điểm $E(3;1;1)$. Khoảng cách (mét) từ đèn đến tấm kính bằng bao nhiêu (làm tròn kết quả đến hàng phần trăm)?
\hinh{\begin{tikzpicture}[scale=0.85,font=\footnotesize,line join=round,line cap=round,>=Stealth]
\draw[->,gray] (0.000,0.000)--(-1.260,-0.945) node[below left]{$x$};
\draw[->,gray] (0.000,0.000)--(4.300,0.000) node[right]{$y$};
\draw[->,gray] (0.000,0.000)--(0.000,3.200) node[above]{$z$};
\node[below right,gray] at (0.000,0.000){$O$};
\fill[gray!25,opacity=.8] (-0.936,-0.702)--(3.400,0.000)--(0.000,2.500)--cycle; \draw[thick] (-0.936,-0.702) node[left]{$A$}--(3.400,0.000) node[below]{$B$}--(0.000,2.500) node[left]{$C$}--cycle;
\shade[ball color=gray!20] (0.640,0.030) circle (.1); \node[right] at (0.640,0.030){$E$};
\end{tikzpicture}}

\begin{loigiai}
\giai Mặt phẳng $(ABC)\colon \dfrac{x}{5}+\dfrac{y}{4}+\dfrac{z}{4}=1\Leftrightarrow 4x+5y+5z-20=0$. $d=\dfrac{|8|}{\sqrt{1056}}\approx 0{,}25$. Đáp số: $0{,}25$.
\end{loigiai}
\end{vidu}
\begin{viduds}
Trong một hệ trục $Oxyz$ (đơn vị: m), một tấm kính phẳng được đặt đi qua ba điểm $A(6;0;0)$, $B(0;6;0)$, $C(0;0;3)$. Xét tính đúng sai của các mệnh đề sau:
\hinh{\begin{tikzpicture}[scale=0.85,font=\footnotesize,line join=round,line cap=round,>=Stealth]
\draw[->,gray] (0.000,0.000)--(-1.260,-0.945) node[below left]{$x$};
\draw[->,gray] (0.000,0.000)--(4.300,0.000) node[right]{$y$};
\draw[->,gray] (0.000,0.000)--(0.000,3.200) node[above]{$z$};
\node[below right,gray] at (0.000,0.000){$O$};
\fill[gray!25,opacity=.8] (-0.936,-0.702)--(3.400,0.000)--(0.000,2.500)--cycle; \draw[thick] (-0.936,-0.702) node[left]{$A$}--(3.400,0.000) node[below]{$B$}--(0.000,2.500) node[left]{$C$}--cycle;
\shade[ball color=gray!20] (0.640,0.030) circle (.1); \node[right] at (0.640,0.030){$E$};
\end{tikzpicture}}

\begin{dungsai}
  \y Khoảng cách từ điểm $E(3;5;0)$ đến tấm kính bằng $36$ m. \oDS
  \y Điểm $M(3;3;0)$ nằm trên tấm kính. \oDS
  \y Mặt phẳng chứa tấm kính có một vectơ pháp tuyến là $(1;1;2)$. \oDS
  \y Tấm kính nằm trên mặt phẳng $x+y+2z+6=0$. \oDS
\end{dungsai}
\begin{loigiai}
\giai Phương trình đoạn chắn $\dfrac{x}{6}+\dfrac{y}{6}+\dfrac{z}{3}=1\Leftrightarrow x+y+2z-6=0$. a) Sai: phải chia cho độ dài vectơ pháp tuyến. b) Đúng. c) Đúng. d) Sai: phương trình đúng là $x+y+2z-6=0$.
\end{loigiai}
\end{viduds}

%% ------------------------------------------------ BÀI TẬP
\luyentap
\phanbt{Trắc nghiệm nhiều phương án lựa chọn}
\begin{bai}Trong không gian $Oxyz$, phương trình của mặt phẳng toạ độ $(Oxy)$ là
\pa{$x+y+z=0$}{$y=0$}{$z=0$}{$x=0$}\dapso{C. }\end{bai}
\begin{bai}Trong không gian $Oxyz$, mặt phẳng đi qua điểm $M(0;0;2)$ và có vectơ pháp tuyến $\vec n=(1;1;-2)$ có phương trình là
\pa{$x-y-2z+4=0$}{$x+y-2z-4=0$}{$z+2=0$}{$x+y-2z+4=0$}\dapso{D. $(x-x_0)+\dots=0\Leftrightarrow x+y-2z+4=0$.}\end{bai}
\begin{bai}Trong không gian $Oxyz$, điểm nào dưới đây thuộc mặt phẳng $(P)\colon 2x+y+z-5=0$?
\pa{$Q(0;1;2)$}{$Q(0;-1;2)$}{$Q(2;1;2)$}{$Q(1;0;3)$}\dapso{D. Thay toạ độ vào phương trình $(P)$.}\end{bai}
\begin{bai}Trong không gian $Oxyz$, mặt phẳng $(P)\colon 3x-3y-2z+1=0$ có một vectơ pháp tuyến là
\pa{$\vec n=(3;-3;1)$}{$\vec n=(3;-3;-2)$}{$\vec n=(-3;-2;1)$}{$\vec n=(3;3;-2)$}\dapso{B. $(P)$ có vectơ pháp tuyến $\vec n=(3;-3;-2)$.}\end{bai}
\begin{bai}Trong không gian $Oxyz$, mặt phẳng đi qua điểm $M(-2;-2;0)$ và song song với mặt phẳng $(Q)\colon 2x+3y-z+4=0$ có phương trình là
\pa{$2x+3y-z-10=0$}{$2x+3y-z+11=0$}{$3x+2y-z=0$}{$2x+3y-z+10=0$}\dapso{D. Mặt phẳng cần tìm có vectơ pháp tuyến $(2;3;-1)$: $2x+3y-z+10=0$.}\end{bai}
\begin{bai}Trong không gian $Oxyz$, cho hai điểm $A(3;1;3)$ và $B(1;-2;1)$. Mặt phẳng trung trực của đoạn thẳng $AB$ có phương trình là
\pa{$4x-y+4z-\dfrac{33}{2}=0$}{$2x+3y+2z-\dfrac{13}{2}=0$}{$2x+3y+2z+2=0$}{$2x+3y+2z-15=0$}\dapso{B. Đi qua trung điểm $I(2;- \dfrac{1}{2};2)$, vectơ pháp tuyến $\vt{AB}=(-2;-3;-2)$: $2x+3y+2z-\dfrac{13}{2}=0$.}\end{bai}
\begin{bai}Trong không gian $Oxyz$, cho ba điểm $A(-1;1;3)$, $B(-2;-1;2)$, $C(2;-1;-2)$. Phương trình mặt phẳng $(ABC)$ là
\pa{$9x-8y+8z-7=0$}{$x-y+z-1=0$}{$x+2y+z-4=0$}{$x-y+z=0$}\dapso{B. $\vt{AB}=(-1;-2;-1)$, $\vt{AC}=(3;-2;-5)$, $\vec n=[\vt{AB},\vt{AC}]=(8;-8;8)$; $(ABC)\colon x-y+z-1=0$.}\end{bai}
\begin{bai}Trong không gian $Oxyz$, khoảng cách từ điểm $M(1;-2;1)$ đến mặt phẳng $(P)\colon x+y-2z+7=0$ bằng
\pa{$4$}{$\dfrac{2 \sqrt{6}}{3}$}{$1 + \dfrac{2 \sqrt{6}}{3}$}{$\dfrac{2}{3}$}\dapso{B. $\mathrm d(M,(P))=\dfrac{4}{\sqrt{6}}=\dfrac{2 \sqrt{6}}{3}$.}\end{bai}
\begin{bai}Trong không gian $Oxyz$, khoảng cách giữa hai mặt phẳng song song $(P)\colon x-2y+2z+4=0$ và $(Q)\colon x-2y+2z-5=0$ bằng
\pa{$\dfrac{1}{3}$}{$3$}{$4$}{$9$}\dapso{B. Lấy một điểm thuộc $(P)$ rồi tính khoảng cách đến $(Q)$: $d=\dfrac{|4-(-5)|}{3}=3$.}\end{bai}
\begin{bai}Trong không gian $Oxyz$, mặt phẳng $(P)\colon \dfrac x1+\dfrac y2+\dfrac z4=1$ cắt các trục $Ox$, $Oy$, $Oz$ lần lượt tại $A$, $B$, $C$. Thể tích khối tứ diện $OABC$ bằng
\pa{$\dfrac{8}{3}$}{$8$}{$\dfrac{4}{3}$}{$4$}\dapso{C. $V=\dfrac16\cdot1\cdot2\cdot4=\dfrac{4}{3}$.}\end{bai}
\begin{bai}Trong một hệ trục $Oxyz$ (đơn vị: mét), một tấm kính phẳng được đặt đi qua ba điểm $A(8;0;0)$, $B(0;6;0)$, $C(0;0;3)$. Một đèn nhỏ đặt tại điểm $E(7;4;1)$. Khoảng cách (mét) từ đèn đến tấm kính (làm tròn đến hàng phần trăm) là
\hinh{\begin{tikzpicture}[scale=0.85,font=\footnotesize,line join=round,line cap=round,>=Stealth]
\draw[->,gray] (0.000,0.000)--(-1.260,-0.945) node[below left]{$x$};
\draw[->,gray] (0.000,0.000)--(4.300,0.000) node[right]{$y$};
\draw[->,gray] (0.000,0.000)--(0.000,3.200) node[above]{$z$};
\node[below right,gray] at (0.000,0.000){$O$};
\fill[gray!25,opacity=.8] (-0.936,-0.702)--(3.400,0.000)--(0.000,2.500)--cycle; \draw[thick] (-0.936,-0.702) node[left]{$A$}--(3.400,0.000) node[below]{$B$}--(0.000,2.500) node[left]{$C$}--cycle;
\shade[ball color=gray!20] (0.640,0.030) circle (.1); \node[right] at (0.640,0.030){$E$};
\end{tikzpicture}}

\pa{$2{,}73$ m}{$4{,}45$ m}{$2{,}23$ m}{$7{,}31$ m}
\dapso{C. Mặt phẳng $(ABC)\colon \dfrac{x}{8}+\dfrac{y}{6}+\dfrac{z}{3}=1\Leftrightarrow 3x+4y+8z-24=0$. $d=\dfrac{|126|}{\sqrt{3204}}\approx 2{,}23$.}\end{bai}
\begin{bai}Trong không gian $Oxyz$ (đơn vị: mét), mặt một sườn đồi được mô hình hoá bởi mặt phẳng $(P)\colon -6x+6y+7z=0$. Một flycam đang ở vị trí $F(-6;6;6)$. Khoảng cách (mét) từ flycam đến mặt sườn đồi là
\hinh{\begin{tikzpicture}[scale=0.85,font=\footnotesize,line join=round,line cap=round,>=Stealth]
\draw[->,gray] (0.000,0.000)--(-1.440,-1.080) node[below left]{$x$};
\draw[->,gray] (0.000,0.000)--(5.000,0.000) node[right]{$y$};
\draw[->,gray] (0.000,0.000)--(0.000,3.300) node[above]{$z$};
\node[below right,gray] at (0.000,0.000){$O$};
\fill[gray!10] (-0.500,0.000)--(-1.760,-0.945)--(3.240,-0.945)--(4.500,0.000)--cycle;
\draw[dashed] (0.300,1.525)--(0.300,-0.675); \fill (0.300,-0.675) circle (.04);
\filldraw[fill=gray!40] (0.300,1.525) ++(-.18,0) rectangle ++(.36,.1); \node[above] at (0.300,1.525){$A$};
\draw[dashed] (3.240,2.530)--(3.240,-0.270); \fill (3.240,-0.270) circle (.04);
\filldraw[fill=gray!40] (3.240,2.530) ++(-.18,0) rectangle ++(.36,.1); \node[above] at (3.240,2.530){$B$};
\draw[thick,dotted] (0.300,1.525)--(3.240,2.530);
\end{tikzpicture}}

\pa{$\dfrac{228}{11}$ m}{$\dfrac{125}{11}$ m}{$114$ m}{$\dfrac{114}{11}$ m}
\dapso{D. $d(F,(P))=\dfrac{|114|}{11}=\dfrac{114}{11}$ (m).}\end{bai}
\begin{bai}Trong một hệ trục $Oxyz$ (đơn vị: mét), một tấm kính phẳng được đặt đi qua ba điểm $A(4;0;0)$, $B(0;5;0)$, $C(0;0;4)$. Một đèn nhỏ đặt tại điểm $E(1;4;2)$. Khoảng cách (mét) từ đèn đến tấm kính (làm tròn đến hàng phần trăm) là
\hinh{\begin{tikzpicture}[scale=0.85,font=\footnotesize,line join=round,line cap=round,>=Stealth]
\draw[->,gray] (0.000,0.000)--(-1.260,-0.945) node[below left]{$x$};
\draw[->,gray] (0.000,0.000)--(4.300,0.000) node[right]{$y$};
\draw[->,gray] (0.000,0.000)--(0.000,3.200) node[above]{$z$};
\node[below right,gray] at (0.000,0.000){$O$};
\fill[gray!25,opacity=.8] (-0.936,-0.702)--(3.400,0.000)--(0.000,2.500)--cycle; \draw[thick] (-0.936,-0.702) node[left]{$A$}--(3.400,0.000) node[below]{$B$}--(0.000,2.500) node[left]{$C$}--cycle;
\shade[ball color=gray!20] (0.640,0.030) circle (.1); \node[right] at (0.640,0.030){$E$};
\end{tikzpicture}}

\pa{$2{,}71$ m}{$6{,}28$ m}{$1{,}85$ m}{$1{,}35$ m}
\dapso{D. Mặt phẳng $(ABC)\colon \dfrac{x}{4}+\dfrac{y}{5}+\dfrac{z}{4}=1\Leftrightarrow 5x+4y+5z-20=0$. $d=\dfrac{|44|}{\sqrt{1056}}\approx 1{,}35$.}\end{bai}
\begin{bai}Trong một hệ trục $Oxyz$ (đơn vị: mét), một tấm kính phẳng được đặt đi qua ba điểm $A(8;0;0)$, $B(0;6;0)$, $C(0;0;4)$. Một đèn nhỏ đặt tại điểm $E(3;2;1)$. Khoảng cách (mét) từ đèn đến tấm kính (làm tròn đến hàng phần trăm) là
\hinh{\begin{tikzpicture}[scale=0.85,font=\footnotesize,line join=round,line cap=round,>=Stealth]
\draw[->,gray] (0.000,0.000)--(-1.260,-0.945) node[below left]{$x$};
\draw[->,gray] (0.000,0.000)--(4.300,0.000) node[right]{$y$};
\draw[->,gray] (0.000,0.000)--(0.000,3.200) node[above]{$z$};
\node[below right,gray] at (0.000,0.000){$O$};
\fill[gray!25,opacity=.8] (-0.936,-0.702)--(3.400,0.000)--(0.000,2.500)--cycle; \draw[thick] (-0.936,-0.702) node[left]{$A$}--(3.400,0.000) node[below]{$B$}--(0.000,2.500) node[left]{$C$}--cycle;
\shade[ball color=gray!20] (0.640,0.030) circle (.1); \node[right] at (0.640,0.030){$E$};
\end{tikzpicture}}

\pa{$0{,}63$ m}{$0{,}13$ m}{$0{,}02$ m}{$0{,}26$ m}
\dapso{B. Mặt phẳng $(ABC)\colon \dfrac{x}{8}+\dfrac{y}{6}+\dfrac{z}{4}=1\Leftrightarrow 3x+4y+6z-24=0$. $d=\dfrac{|-8|}{\sqrt{3904}}\approx 0{,}13$.}\end{bai}

\phanbt{Trắc nghiệm đúng/sai}
\begin{bai}Trong không gian $Oxyz$, cho mặt phẳng $(P)\colon x+2y-2z+5=0$.
\begin{dungsai}
  \y Vectơ $\vec n=(2;4;-4)$ là một vectơ pháp tuyến của $(P)$. \oDS
  \y Điểm $A(0;0;\dfrac{5}{2})$ thuộc $(P)$. \oDS
  \y Khoảng cách từ điểm $M(1;1;1)$ đến $(P)$ bằng $2$. \oDS
  \y Mặt phẳng $(Q)\colon 2x+y-2z-1=0$ vuông góc với $(P)$. \oDS
\end{dungsai}
\dapso{a) Đ. b) Đ. c) Đ. d) S. $(2;4;-4)=2(1;2;-2)$; $d(M,(P))=2$; $(1;2;-2)\cdot(2;1;2)=0$, còn $(1;2;-2)\cdot(2;1;-2)=8\ne0$.}\end{bai}
\begin{bai}Trong không gian $Oxyz$, cho ba điểm $A(1;0;0)$, $B(0;2;0)$, $C(0;0;3)$.
\begin{dungsai}
  \y Mặt phẳng $(ABC)$ có phương trình $\dfrac x1+\dfrac y2+\dfrac z3=1$. \oDS
  \y Vectơ $\vec n=(6;3;2)$ là một vectơ pháp tuyến của $(ABC)$. \oDS
  \y Khoảng cách từ gốc toạ độ $O$ đến $(ABC)$ bằng $\dfrac67$. \oDS
  \y Thể tích khối tứ diện $OABC$ bằng $1$. \oDS
\end{dungsai}
\dapso{a) Đ. b) Đ. c) Đ. d) Đ. $(ABC)\colon 6x+3y+2z-6=0$; $d(O,(ABC))=\dfrac{6}{7}$; $V=\dfrac{1\cdot2\cdot3}{6}=1$.}\end{bai}
\begin{bai}Trong một hệ trục $Oxyz$ (đơn vị: m), một tấm kính phẳng được đặt đi qua ba điểm $A(8;0;0)$, $B(0;3;0)$, $C(0;0;4)$. Xét tính đúng sai của các mệnh đề sau:
\hinh{\begin{tikzpicture}[scale=0.85,font=\footnotesize,line join=round,line cap=round,>=Stealth]
\draw[->,gray] (0.000,0.000)--(-1.260,-0.945) node[below left]{$x$};
\draw[->,gray] (0.000,0.000)--(4.300,0.000) node[right]{$y$};
\draw[->,gray] (0.000,0.000)--(0.000,3.200) node[above]{$z$};
\node[below right,gray] at (0.000,0.000){$O$};
\fill[gray!25,opacity=.8] (-0.936,-0.702)--(3.400,0.000)--(0.000,2.500)--cycle; \draw[thick] (-0.936,-0.702) node[left]{$A$}--(3.400,0.000) node[below]{$B$}--(0.000,2.500) node[left]{$C$}--cycle;
\shade[ball color=gray!20] (0.640,0.030) circle (.1); \node[right] at (0.640,0.030){$E$};
\end{tikzpicture}}

\begin{dungsai}
  \y Tấm kính nằm trên mặt phẳng $3x+8y+6z-24=0$. \oDS
  \y Khoảng cách từ điểm $E(6;2;1)$ đến tấm kính bằng $64$ m. \oDS
  \y Tấm kính nằm trên mặt phẳng $3x+8y+6z+24=0$. \oDS
  \y Điểm $M(4;\dfrac{3}{2};0)$ nằm trên tấm kính. \oDS
\end{dungsai}
\dapso{Phương trình đoạn chắn $\dfrac{x}{8}+\dfrac{y}{3}+\dfrac{z}{4}=1\Leftrightarrow 3x+8y+6z-24=0$. a) Đúng. b) Sai: phải chia cho độ dài vectơ pháp tuyến. c) Sai: phương trình đúng là $3x+8y+6z-24=0$. d) Đúng.}\end{bai}
\begin{bai}Trong một hệ trục $Oxyz$ (đơn vị: m), một tấm kính phẳng được đặt đi qua ba điểm $A(6;0;0)$, $B(0;4;0)$, $C(0;0;4)$. Xét tính đúng sai của các mệnh đề sau:
\hinh{\begin{tikzpicture}[scale=0.85,font=\footnotesize,line join=round,line cap=round,>=Stealth]
\draw[->,gray] (0.000,0.000)--(-1.260,-0.945) node[below left]{$x$};
\draw[->,gray] (0.000,0.000)--(4.300,0.000) node[right]{$y$};
\draw[->,gray] (0.000,0.000)--(0.000,3.200) node[above]{$z$};
\node[below right,gray] at (0.000,0.000){$O$};
\fill[gray!25,opacity=.8] (-0.936,-0.702)--(3.400,0.000)--(0.000,2.500)--cycle; \draw[thick] (-0.936,-0.702) node[left]{$A$}--(3.400,0.000) node[below]{$B$}--(0.000,2.500) node[left]{$C$}--cycle;
\shade[ball color=gray!20] (0.640,0.030) circle (.1); \node[right] at (0.640,0.030){$E$};
\end{tikzpicture}}

\begin{dungsai}
  \y Tấm kính nằm trên mặt phẳng $2x+3y+3z-12=0$. \oDS
  \y Khoảng cách từ điểm $E(3;3;1)$ đến tấm kính bằng $\dfrac{3 \sqrt{22}}{11}$ m. \oDS
  \y Khoảng cách từ điểm $E(3;3;1)$ đến tấm kính bằng $48$ m. \oDS
  \y Tấm kính nằm trên mặt phẳng $2x+3y+3z+12=0$. \oDS
\end{dungsai}
\dapso{Phương trình đoạn chắn $\dfrac{x}{6}+\dfrac{y}{4}+\dfrac{z}{4}=1\Leftrightarrow 2x+3y+3z-12=0$. a) Đúng. b) Đúng. c) Sai: phải chia cho độ dài vectơ pháp tuyến. d) Sai: phương trình đúng là $2x+3y+3z-12=0$.}\end{bai}

\phanbt{Trả lời ngắn}
\begin{bai}Trong không gian $Oxyz$, mặt phẳng đi qua ba điểm $A(3;0;-2)$, $B(1;2;-1)$, $C(0;-2;3)$ có phương trình dạng $ax+by+cz+d=0$ với $a$, $b$, $c$, $d$ là các số nguyên tối giản (ước chung lớn nhất bằng $1$) và $a>0$ (nếu $a=0$ thì xét hệ số khác $0$ đầu tiên dương). Tính $a+b+c+d$.\otraloi\dapso{13. $\vec n=(12;7;10)$; $12x+7y+10z-16=0$; tổng bằng $13$.}\end{bai}
\begin{bai}Trong không gian $Oxyz$, tính khoảng cách từ điểm $M(-2;2;0)$ đến mặt phẳng $(P)\colon 2x+2y+z-4=0$ (làm tròn đến hàng phần trăm).\otraloi\dapso{1{,}33. $d=\dfrac{4}{3}\approx1{,}33$.}\end{bai}
\begin{bai}Trong không gian $Oxyz$, mặt phẳng đi qua ba điểm $A(-2;3;0)$, $B(2;-2;-1)$, $C(3;2;0)$ có phương trình dạng $ax+by+cz+d=0$ với $a$, $b$, $c$, $d$ là các số nguyên tối giản (ước chung lớn nhất bằng $1$) và $a>0$ (nếu $a=0$ thì xét hệ số khác $0$ đầu tiên dương). Tính $a+b+c+d$.\otraloi\dapso{-28. $\vec n=(-1;-5;21)$; $x+5y-21z-13=0$; tổng bằng $-28$.}\end{bai}
\begin{bai}Trong một căn phòng, chọn hệ trục $Oxyz$ (đơn vị mét) sao cho sàn nhà thuộc mặt phẳng $(Oxy)$ và một bức tường nằm trong mặt phẳng $(P)\colon 3x+4y-24=0$. Một chiếc đèn được treo tại điểm $D(1;2;5)$. Tính khoảng cách từ đèn đến bức tường (đơn vị mét).\otraloi\dapso{2{,}6. $d=\dfrac{|3+8-24|}{5}=2{,}6$.}\end{bai}
\begin{bai}Trong một căn phòng, chọn hệ trục $Oxyz$ (đơn vị mét) sao cho sàn nhà thuộc mặt phẳng $(Oxy)$ và một bức tường nằm trong mặt phẳng $(P)\colon 3x+4y-24=0$. Một chiếc đèn được treo tại điểm $D(1;2;3)$. Tính khoảng cách từ đèn đến bức tường (đơn vị mét).\otraloi\dapso{2{,}6. $d=\dfrac{|3+8-24|}{5}=2{,}6$.}\end{bai}
\begin{bai}Trong không gian $Oxyz$, mặt phẳng $(P)$ đi qua ba điểm $A(3;0;0)$, $B(0;4;0)$, $C(0;0;6)$. Tính thể tích khối tứ diện $OABC$.\otraloi\dapso{12. $V=\dfrac{3\cdot4\cdot6}{6}=12$.}\end{bai}
\begin{bai}Trong không gian $Oxyz$ (đơn vị: mét), mặt một sườn đồi được mô hình hoá bởi mặt phẳng $(P)\colon z+7=0$. Một flycam đang ở vị trí $F(-1;3;7)$. Khoảng cách (mét) từ flycam đến mặt sườn đồi bằng bao nhiêu?
\hinh{\begin{tikzpicture}[scale=0.85,font=\footnotesize,line join=round,line cap=round,>=Stealth]
\draw[->,gray] (0.000,0.000)--(-1.440,-1.080) node[below left]{$x$};
\draw[->,gray] (0.000,0.000)--(5.000,0.000) node[right]{$y$};
\draw[->,gray] (0.000,0.000)--(0.000,3.300) node[above]{$z$};
\node[below right,gray] at (0.000,0.000){$O$};
\fill[gray!10] (-0.500,0.000)--(-1.760,-0.945)--(3.240,-0.945)--(4.500,0.000)--cycle;
\draw[dashed] (0.300,1.525)--(0.300,-0.675); \fill (0.300,-0.675) circle (.04);
\filldraw[fill=gray!40] (0.300,1.525) ++(-.18,0) rectangle ++(.36,.1); \node[above] at (0.300,1.525){$A$};
\draw[dashed] (3.240,2.530)--(3.240,-0.270); \fill (3.240,-0.270) circle (.04);
\filldraw[fill=gray!40] (3.240,2.530) ++(-.18,0) rectangle ++(.36,.1); \node[above] at (3.240,2.530){$B$};
\draw[thick,dotted] (0.300,1.525)--(3.240,2.530);
\end{tikzpicture}}
\otraloi
\dapso{14. $d(F,(P))=\dfrac{|14|}{1}=14$ (m).}\end{bai}
\begin{bai}Trong một hệ trục $Oxyz$ (đơn vị: mét), một tấm kính phẳng được đặt đi qua ba điểm $A(7;0;0)$, $B(0;6;0)$, $C(0;0;4)$. Một đèn nhỏ đặt tại điểm $E(3;1;3)$. Khoảng cách (mét) từ đèn đến tấm kính bằng bao nhiêu (làm tròn kết quả đến hàng phần trăm)?
\hinh{\begin{tikzpicture}[scale=0.85,font=\footnotesize,line join=round,line cap=round,>=Stealth]
\draw[->,gray] (0.000,0.000)--(-1.260,-0.945) node[below left]{$x$};
\draw[->,gray] (0.000,0.000)--(4.300,0.000) node[right]{$y$};
\draw[->,gray] (0.000,0.000)--(0.000,3.200) node[above]{$z$};
\node[below right,gray] at (0.000,0.000){$O$};
\fill[gray!25,opacity=.8] (-0.936,-0.702)--(3.400,0.000)--(0.000,2.500)--cycle; \draw[thick] (-0.936,-0.702) node[left]{$A$}--(3.400,0.000) node[below]{$B$}--(0.000,2.500) node[left]{$C$}--cycle;
\shade[ball color=gray!20] (0.640,0.030) circle (.1); \node[right] at (0.640,0.030){$E$};
\end{tikzpicture}}
\otraloi
\dapso{1,04. Mặt phẳng $(ABC)\colon \dfrac{x}{7}+\dfrac{y}{6}+\dfrac{z}{4}=1\Leftrightarrow 12x+14y+21z-84=0$. $d=\dfrac{|58|}{\sqrt{3124}}\approx 1{,}04$.}\end{bai}
\begin{bai}Trong không gian $Oxyz$ (đơn vị: mét), mặt một sườn đồi được mô hình hoá bởi mặt phẳng $(P)\colon -4x+3y+12z+3=0$. Một flycam đang ở vị trí $F(5;-3;6)$. Khoảng cách (mét) từ flycam đến mặt sườn đồi bằng bao nhiêu (làm tròn kết quả đến hàng phần trăm)?
\hinh{\begin{tikzpicture}[scale=0.85,font=\footnotesize,line join=round,line cap=round,>=Stealth]
\draw[->,gray] (0.000,0.000)--(-1.440,-1.080) node[below left]{$x$};
\draw[->,gray] (0.000,0.000)--(5.000,0.000) node[right]{$y$};
\draw[->,gray] (0.000,0.000)--(0.000,3.300) node[above]{$z$};
\node[below right,gray] at (0.000,0.000){$O$};
\fill[gray!10] (-0.500,0.000)--(-1.760,-0.945)--(3.240,-0.945)--(4.500,0.000)--cycle;
\draw[dashed] (0.300,1.525)--(0.300,-0.675); \fill (0.300,-0.675) circle (.04);
\filldraw[fill=gray!40] (0.300,1.525) ++(-.18,0) rectangle ++(.36,.1); \node[above] at (0.300,1.525){$A$};
\draw[dashed] (3.240,2.530)--(3.240,-0.270); \fill (3.240,-0.270) circle (.04);
\filldraw[fill=gray!40] (3.240,2.530) ++(-.18,0) rectangle ++(.36,.1); \node[above] at (3.240,2.530){$B$};
\draw[thick,dotted] (0.300,1.525)--(3.240,2.530);
\end{tikzpicture}}
\otraloi
\dapso{3,54. $d(F,(P))=\dfrac{|46|}{13}=\dfrac{46}{13}$ (m).}\end{bai}
\phieutln

\phanbt{Tự luận}
\begin{bai}Trong không gian $Oxyz$, cho hình hộp $ABCD.A'B'C'D'$ với $A(1;-1;3)$, $B(0;2;4)$, $D(2;-1;1)$, $A'(0;1;2)$. a) Tính toạ độ các điểm $C$, $B'$, $D'$. b) Viết phương trình mặt phẳng $(CB'D')$.
\dongke{7}
\dapso{a) $C(1;2;2)$, $B'(-1;4;3)$, $D'(1;1;0)$. b) $\vec n=\big[\vt{CB'},\vt{CD'}\big]=(-3;-4;2)$; $(CB'D')\colon3x+4y-2z-7=0$.}\end{bai}
\begin{bai}Trong không gian $Oxyz$, viết phương trình mặt phẳng $(P)$ đi qua điểm $M(1;-1;5)$ và vuông góc với hai mặt phẳng $(Q)\colon3x+2y-z=0$, $(R)\colon x+y-z=0$.
\dongke{5}
\dapso{$\vec n_P=[\vec n_Q,\vec n_R]=(-1;2;1)$; $(P)\colon x-2y-z+2=0$.}\end{bai}
\begin{bai}Trong không gian $Oxyz$, cho hai mặt phẳng $(P)\colon x+3y-z=0$ và $(Q)\colon x-y-2z+1=0$. a) Chứng minh $(P)\perp(Q)$. b) Tìm điểm $M$ thuộc trục $Ox$ và cách đều hai mặt phẳng $(P)$, $(Q)$.
\dongke{6}
\dapso{a) $1-3+2=0$. b) $M(t;0;0)$: $\dfrac{|t|}{\sqrt{11}}=\dfrac{|t+1|}{\sqrt6}\Leftrightarrow5t^2+22t+11=0\Leftrightarrow t=\dfrac{-11\pm\sqrt{66}}5$.}\end{bai}
\begin{bai}Trong không gian $Oxyz$ (đơn vị: mét), mặt một sườn đồi được mô hình hoá bởi mặt phẳng $(P)\colon -2x+2y+z+3=0$. Một flycam đang ở vị trí $F(-2;5;13)$. Khoảng cách (mét) từ flycam đến mặt sườn đồi bằng bao nhiêu?
\hinh{\begin{tikzpicture}[scale=0.85,font=\footnotesize,line join=round,line cap=round,>=Stealth]
\draw[->,gray] (0.000,0.000)--(-1.440,-1.080) node[below left]{$x$};
\draw[->,gray] (0.000,0.000)--(5.000,0.000) node[right]{$y$};
\draw[->,gray] (0.000,0.000)--(0.000,3.300) node[above]{$z$};
\node[below right,gray] at (0.000,0.000){$O$};
\fill[gray!10] (-0.500,0.000)--(-1.760,-0.945)--(3.240,-0.945)--(4.500,0.000)--cycle;
\draw[dashed] (0.300,1.525)--(0.300,-0.675); \fill (0.300,-0.675) circle (.04);
\filldraw[fill=gray!40] (0.300,1.525) ++(-.18,0) rectangle ++(.36,.1); \node[above] at (0.300,1.525){$A$};
\draw[dashed] (3.240,2.530)--(3.240,-0.270); \fill (3.240,-0.270) circle (.04);
\filldraw[fill=gray!40] (3.240,2.530) ++(-.18,0) rectangle ++(.36,.1); \node[above] at (3.240,2.530){$B$};
\draw[thick,dotted] (0.300,1.525)--(3.240,2.530);
\end{tikzpicture}}

\dongke{6}
\dapso{$d(F,(P))=\dfrac{|30|}{3}=10$ (m).}\end{bai}
\begin{bai}Trong không gian $Oxyz$ (đơn vị: mét), mặt một sườn đồi được mô hình hoá bởi mặt phẳng $(P)\colon 3x+4z+4=0$. Một flycam đang ở vị trí $F(2;-4;15)$. Khoảng cách (mét) từ flycam đến mặt sườn đồi bằng bao nhiêu?
\hinh{\begin{tikzpicture}[scale=0.85,font=\footnotesize,line join=round,line cap=round,>=Stealth]
\draw[->,gray] (0.000,0.000)--(-1.440,-1.080) node[below left]{$x$};
\draw[->,gray] (0.000,0.000)--(5.000,0.000) node[right]{$y$};
\draw[->,gray] (0.000,0.000)--(0.000,3.300) node[above]{$z$};
\node[below right,gray] at (0.000,0.000){$O$};
\fill[gray!10] (-0.500,0.000)--(-1.760,-0.945)--(3.240,-0.945)--(4.500,0.000)--cycle;
\draw[dashed] (0.300,1.525)--(0.300,-0.675); \fill (0.300,-0.675) circle (.04);
\filldraw[fill=gray!40] (0.300,1.525) ++(-.18,0) rectangle ++(.36,.1); \node[above] at (0.300,1.525){$A$};
\draw[dashed] (3.240,2.530)--(3.240,-0.270); \fill (3.240,-0.270) circle (.04);
\filldraw[fill=gray!40] (3.240,2.530) ++(-.18,0) rectangle ++(.36,.1); \node[above] at (3.240,2.530){$B$};
\draw[thick,dotted] (0.300,1.525)--(3.240,2.530);
\end{tikzpicture}}

\dongke{6}
\dapso{$d(F,(P))=\dfrac{|70|}{5}=14$ (m).}\end{bai}

\phanbt{Bài tập tự luận rèn luyện}
\begin{bai}Trong không gian $Oxyz$, cho ba điểm $A(1;0;2)$, $B(2;1;0)$ và $C(0;3;1)$.
\begin{enumerate}[label=\alph*)]
\item Tìm toạ độ các vectơ $\overrightarrow{AB}$, $\overrightarrow{AC}$.
\item Tính tích có hướng $\left[\overrightarrow{AB},\overrightarrow{AC}\right]$.
\item Viết phương trình mặt phẳng $(ABC)$ và kiểm tra lại bằng cách thay toạ độ ba điểm.
\item Tính diện tích tam giác $ABC$.
\item Tính khoảng cách từ gốc toạ độ $O$ đến mặt phẳng $(ABC)$ (làm tròn đến hàng phần trăm).
\item Viết phương trình mặt phẳng đi qua $O$ và song song với mặt phẳng $(ABC)$.
\end{enumerate}
\dapso{a) $(1;1;-2)$, $(-1;3;-1)$. b) $(5;3;4)$. c) $5x+3y+4z-13=0$. d) $\dfrac12\sqrt{50}=\dfrac{5\sqrt2}2$. e) $\dfrac{13}{\sqrt{50}}\approx1{,}84$. f) $5x+3y+4z=0$.}\end{bai}

\begin{bai}Trong không gian $Oxyz$, cho mặt phẳng $(P)\colon 2x-y+2z-3=0$ và điểm $N(3;-2;4)$.
\begin{enumerate}[label=\alph*)]
\item Chỉ ra một vectơ pháp tuyến của $(P)$.
\item Điểm $M(1;1;1)$ có thuộc $(P)$ không?
\item Tính khoảng cách từ $N$ đến mặt phẳng $(P)$.
\item Viết phương trình mặt phẳng $(Q)$ đi qua $N$ và song song với $(P)$.
\item Tính khoảng cách giữa hai mặt phẳng $(P)$ và $(Q)$; so sánh với kết quả câu c) và giải thích.
\item Viết phương trình mặt phẳng đi qua $N$, đồng thời vuông góc với cả $(P)$ và mặt phẳng $(Oxy)$.
\end{enumerate}
\dapso{a) $(2;-1;2)$. b) Có, vì $2-1+2-3=0$. c) $\dfrac{13}3$. d) $2x-y+2z-16=0$. e) $\dfrac{|-16+3|}3=\dfrac{13}3$, bằng câu c) vì $N\in(Q)$ và $(Q)\parallel(P)$. f) $\vec n=[(2;-1;2),(0;0;1)]=(-1;-2;0)$: $x+2y+1=0$.}\end{bai}

\begin{bai}Một mái hiên bằng tôn được lắp nghiêng phía trước một ngôi nhà. Chọn hệ trục $Oxyz$ (đơn vị: mét) có mặt phẳng $(Oxy)$ là mặt sân, trục $Oz$ hướng thẳng đứng lên trên. Mái hiên là một phần mặt phẳng đi qua ba điểm $A(0;0;3)$, $B(8;0;3)$ và $C(0;4;5)$, trong đó cạnh $AB$ là mép dưới và mép trên của mái nằm sát tường nhà.
\hinh{\begin{tikzpicture}[scale=0.85,font=\footnotesize,line join=round,line cap=round,>=Stealth]
\draw[thick] (-.5,0)--(5.6,0); \fill[pattern=north east lines] (-.5,-.15) rectangle (5.6,0);
\filldraw[fill=gray!10] (1,0) rectangle (5,2.6); \node at (3.5,.6){tường nhà};
\filldraw[fill=gray!40] (0,1.5)--(4,1.5)--(5,2.2)--(1,2.2)--cycle; \draw (0,1.5)--(0,0) (4,1.5)--(4,0);
\fill (0,1.5) circle (.04) node[left]{$A$}; \fill (4,1.5) circle (.04) node[below right]{$B$}; \fill (1,2.2) circle (.04) node[above left]{$C$};
\node at (2.5,1.85){mái hiên};
\end{tikzpicture}}
\begin{enumerate}[label=\alph*)]
\item Tìm toạ độ các vectơ $\overrightarrow{AB}$, $\overrightarrow{AC}$.
\item Tìm một vectơ pháp tuyến của mặt phẳng chứa mái hiên.
\item Viết phương trình mặt phẳng chứa mái hiên.
\item Một bóng đèn treo tại điểm $E(4;2;1)$ dưới mái hiên. Tính khoảng cách từ bóng đèn tới mặt phẳng mái (làm tròn đến hàng phần trăm).
\item Mép trên của mái nằm ở độ cao bao nhiêu mét so với mặt sân?
\item Mái hiên có dạng hình chữ nhật với hai cạnh là $AB$ và $AC$. Tính diện tích tôn cần dùng (làm tròn đến hàng phần trăm).
\end{enumerate}
\dapso{a) $(8;0;0)$, $(0;4;2)$. b) $[\overrightarrow{AB},\overrightarrow{AC}]=(0;-16;32)$, chọn $(0;1;-2)$. c) $y-2z+6=0$. d) $\dfrac6{\sqrt5}\approx2{,}68$ m. e) $5$ m. f) $8\cdot2\sqrt5=16\sqrt5\approx35{,}78$ m$^2$.}\end{bai}

\begin{bai}Trong không gian $Oxyz$, cho ba điểm $A(2;0;0)$, $B(0;3;0)$, $C(0;0;6)$ lần lượt nằm trên ba trục toạ độ.
\begin{enumerate}[label=\alph*)]
\item Viết phương trình mặt phẳng $(ABC)$ dưới dạng đoạn chắn, rồi đưa về dạng tổng quát với hệ số nguyên.
\item Tính thể tích khối tứ diện $OABC$.
\item Tính khoảng cách từ $O$ đến mặt phẳng $(ABC)$.
\item Từ câu b) và c), suy ra diện tích tam giác $ABC$.
\item Tìm toạ độ hình chiếu vuông góc $H$ của $O$ trên mặt phẳng $(ABC)$.
\item Tìm toạ độ trọng tâm $G$ của tam giác $ABC$ và kiểm tra xem $H$ có trùng $G$ không.
\end{enumerate}
\dapso{a) $\dfrac x2+\dfrac y3+\dfrac z6=1\Leftrightarrow3x+2y+z-6=0$. b) $6$. c) $\dfrac6{\sqrt{14}}$. d) $S=\dfrac{3V}d=3\sqrt{14}$. e) $H=t(3;2;1)$, $14t=6$: $H\left(\dfrac97;\dfrac67;\dfrac37\right)$. f) $G\left(\dfrac23;1;2\right)$, không trùng $H$.}\end{bai}

\begin{bai}Trong không gian $Oxyz$, cho ba mặt phẳng $(P)\colon x+2y-2z+1=0$, $(Q)\colon 2x+4y-4z+5=0$ và $(R)\colon 2x-y+3=0$.
\begin{enumerate}[label=\alph*)]
\item Chứng minh $(P)$ và $(Q)$ song song với nhau.
\item Tính khoảng cách giữa $(P)$ và $(Q)$.
\item Chứng minh $(P)$ vuông góc với $(R)$.
\item Viết phương trình mặt phẳng song song và cách đều hai mặt phẳng $(P)$, $(Q)$.
\item Điểm $M(1;-1;2)$ có nằm giữa hai mặt phẳng $(P)$ và $(Q)$ không?
\item Viết phương trình các mặt phẳng song song với $(P)$ và cách $(P)$ một khoảng bằng $2$.
\end{enumerate}
\dapso{a) $\dfrac12=\dfrac24=\dfrac{-2}{-4}\ne\dfrac15$. b) $(Q)\colon x+2y-2z+\dfrac52=0$; khoảng cách $\dfrac12$. c) $\vec n_P\cdot\vec n_R=2-2+0=0$. d) $x+2y-2z+\dfrac74=0$. e) Thay $M$ vào hai vế trái (cùng dạng $x+2y-2z+d$) được $-4$ và $-2{,}5$ cùng dấu: $M$ không nằm giữa. f) $x+2y-2z+7=0$ và $x+2y-2z-5=0$.}\end{bai}
